\section{Related Work}
\label{sec:related}

\para{Self-evolving agents and skill libraries.} A growing literature builds agents that
accumulate reusable procedures or rewrite their own prompts and code:
Voyager~\citep{voyager} maintains an executable skill library, Promptbreeder~\citep{promptbreeder}
evolves task- and mutation-prompts, the Darwin G\"odel Machine~\citep{dgm} rewrites its own
code, and the self-evolving-agents survey~\citep{selfevolvesurvey} formalizes the loop of
inputs, agent, environment, and optimizer. The G\"odel Machine~\citep{godelmachine} is the
theoretical ideal these approximate. A recurring observation, made by Voyager's own authors,
is that such systems harness knowledge already encapsulated in the pretrained model and
perform an in-context form of novelty search. Unlike this line of work, we do not propose a
new self-evolving system; we build a substrate and a discriminator that measure whether such
systems create knowledge or re-elicit it.

\para{Controlled skill ablations.} The closest works measure task performance with and without
a skill while holding the base model fixed. SkillsBench~\citep{skillsbench} reports that
human-curated skills raise pass rate by $+16.6$ points across 18 agent$\times$model
configurations, whereas agent self-generated skills fall $8$--$12$ points \emph{below}
baseline. A data-science skill ablation~\citep{dsskills} finds the full skill condition
statistically indistinguishable from a token-matched irrelevant control (all $p\geq0.396$), and
Library Drift~\citep{librarydrift} shows self-curated skill banks silently degrade without
outcome-driven governance, with a no-injection floor of $+0.002$. Meta-agent
studies~\citep{metaagents} similarly find ``evolved'' designs mostly reshuffle known building
blocks. Building on these, we add the missing control: an explicit test of whether the model
could self-generate the skill, and a substrate where the correct answer is known.

\para{Elicitation versus new capability.} A parallel literature questions whether apparent
gains reflect new capability. Reasoning-or-Reciting~\citep{reasoningreciting} shows that on
counterfactual task variants (\eg base-9 arithmetic) performance collapses while comprehension
stays high, indicating recitation rather than reasoning. The emergent-abilities-as-mirage
critique~\citep{mirage} warns that sharp ``new'' abilities are often metric artifacts.
Crucially, LLMs-cannot-self-correct-yet~\citep{cannotselfcorrect} shows that without an
external signal, self-refinement does not add correctness and often degrades it, and
STaR~\citep{star} works only because an external correctness filter gates self-generated
rationales. Our verifier-only versus with-answer self-evolution experiment directly
instantiates this distinction on our substrate.

\para{Novelty measurement and verifiable discovery.} Idea-novelty studies find LLMs can
produce novel-\emph{sounding} ideas~\citep{novelideas} whose novelty scales with retrieved
external knowledge~\citep{nova}, but HindSight~\citep{hindsight} shows judged novelty
anti-correlates with real impact --- motivating our deterministic verifiers and our refusal to
use LLM-as-judge for primary signals. The one clear case of verifiable novelty,
AlphaEvolve~\citep{alphaevolve}, obtains genuinely new mathematical results from evolutionary
search plus an external ground-truth evaluator, not from the model introspecting; the Red Queen
G\"odel Machine~\citep{redqueen} warns that learning the evaluator invites verifier-hacking.
Our results reproduce and mechanize this pattern: genuine novelty enters a skill only from an
external source, and self-evolution alone cannot supply it.
